\documentclass[11pt]{amsart}
\usepackage{geometry}                % See geometry.pdf to learn the layout options. There are lots.
\geometry{letterpaper}                   % ... or a4paper or a5paper or ... 
%\geometry{landscape}                % Activate for for rotated page geometry
%\usepackage[parfill]{parskip}    % Activate to begin paragraphs with an empty line rather than an indent
\usepackage{graphicx}
\usepackage{amssymb}
\usepackage{epstopdf}
\DeclareGraphicsRule{.tif}{png}{.png}{`convert #1 `dirname #1`/`basename #1 .tif`.png}

\title{PS 3.6}
\author{Enrique Zuñiga Zuñiga}
%\date{}                                           % Activate to display a given date or no date

\begin{document}
\maketitle
%\section{}
\noindent Dado el sueldo de N trabajadores, considere un aumento del 15\% a cada uno de ellos si su sueldo es inferior a \$800. Imprima el sueldo con el aumento incorporado (si corresponde). Haga el diagrama de flujo correspondiente.\newline
Datos:$ N, SUE_1, SUE_2, .. SUE_N$\newline
Donde:\newline
N es una variable de tipo entero que representa el número de empleados de la empresa.\newline
$SUE_i$ es una variable de tipo real que representa el sueldo del trabajador i $(1 < i < N)$.\newline
%\subsection{}



\end{document}  